\section{Experiments}
\label{sec:experiments}

In this section, we present our empirical evaluation of NETA against foundational and state-of-the-art baselines. We detail our experimental setup, present our quantitative findings, and conduct a detailed analysis aligning with the principles of Occam's razor.

\subsection{Experimental Setup}
\textbf{Model Backbone.} Following standard literature in model merging, we employ the pre-trained CLIP ViT-B/32 \cite{radford2021clip} as our model backbone. The model's parameter weights are partitioned into $L = 13$ discrete parameter groups: 12 Transformer blocks (comprising self-attention and MLP layers) and 1 final visual projection head (`model.visual.proj`). 

\textbf{Task Experts and Datasets.} We evaluate multi-task merging performance across four distinct visual classification datasets representing various image domains and complexities: \textbf{MNIST}, \textbf{FashionMNIST}, \textbf{CIFAR-10}, and \textbf{SVHN}. We utilize independently fine-tuned expert checkpoints and their corresponding classification heads downloaded from the public Hugging Face hub repository.

\textbf{Baselines.} We compare our proposed NETA against several prominent baselines representing both zero-shot and test-time optimization paradigms:
\begin{itemize}
    \item \textbf{Task Arithmetic (TA)} \cite{ilharco2023editing}: The foundational zero-shot merging baseline, using a flat global scaling coefficient $\lambda_0 = 0.3$ across all layers and tasks.
    \item \textbf{TIES-Merging} \cite{yadav2023resolving}: A prominent zero-shot merging baseline that prunes parameter updates (keeping top 80\% parameters based on magnitude), resolves sign conflicts via a parameter-wise majority sign election, and averages non-conflicting task updates.
    \item \textbf{DARE} \cite{yu2024dare}: A standard zero-shot merging method that randomly sparsifies task vectors (keeping $90\%$ of updates, i.e., dropping $p_{\text{drop}} = 0.1$) and rescales the remaining updates by $1/(1-p_{\text{drop}})$ to mitigate multi-task representation interference.
    \item \textbf{Task-Wise AdaMerging (TTA)} \cite{yang2024adamerging}: A Test-Time Adaptation baseline that optimizes $K = 4$ task-wise scaling coefficients to minimize prediction entropy over 256 unlabeled calibration images. We optimize using Adam with a learning rate of $5 \times 10^{-3}$ and a batch size of 32 images per task for 20 epochs.
    \item \textbf{Layer-Wise AdaMerging (TTA)} \cite{yang2024adamerging}: The state-of-the-art test-time adaptation method that optimizes $L \times K = 52$ independent layer-wise, task-wise scaling parameters under the same prediction entropy minimization objective, utilizing the same optimization settings and batch size of 32.
\end{itemize}

\textbf{Omitted Baselines.} We explicitly justify the omission of **RegMean** \cite{jin2023regmean} from our quantitative comparison. RegMean is a closed-form baseline discussed in literature, but it requires computing activation covariance matrices on calibration data. This demands access to a dedicated validation set with substantial forward passes, which conflicts with our strict zero-shot, completely data-free evaluation paradigm.

To ensure a fully controlled, apples-to-apples comparison of weight-space transformations under a fixed scaling budget, both TIES-Merging and DARE are evaluated using the same global scaling factor $\lambda_0 = 0.3$ as standard Task Arithmetic. Setting $\lambda_0 = 0.3$ serves as a standard, stable default, although we acknowledge that further hyperparameter tuning of their global coefficients could impact their absolute average performance.

\textbf{Evaluation Protocol and Sub-Sampling.} To ensure scientific rigor and manage the computational overhead of running multiple baselines, optimization sweeps, and seeds under resource constraints, all evaluations are performed on a representative subset of 1024 randomly sampled test images from each dataset. We conduct all experiments across three independent random seeds (42, 100, and 2026) and report the mean multi-task visual classification accuracy (\%) along with its standard deviation. While sub-sampling represents a practical and necessary constraint, using 1024 images provides a robust statistical baseline (where each individual image corresponds to approximately $0.098\%$ accuracy) that enables a highly reliable and controlled comparative analysis of test-time optimization loops.

\begin{table*}[t]
\caption{Multi-task visual classification accuracy (mean \% $\pm$ standard deviation across 3 seeds) for CLIP ViT-B/32 merging. Under zero-shot settings, NETA improves performance over Task Arithmetic, TIES-Merging, and DARE across MNIST and FashionMNIST. Under Test-Time Adaptation, Task-Wise AdaMerging exhibits catastrophic degradation on FashionMNIST due to the Overfitting-Optimizer Paradox, whereas NETA achieves remarkable robustness.}
\label{tab:results}
\vskip 0.15in
\begin{center}
\begin{small}
\begin{sc}
\setlength{\tabcolsep}{4pt}
\begin{tabular}{lcccccr}
\toprule
Method & Opt. Params & MNIST (\%) & FashionMNIST (\%) & CIFAR-10 (\%) & SVHN (\%) & Average (\%) \\
\midrule
\textit{Zero-Shot (Training-Free):} \\
Task Arithmetic & 0 & 96.03 $\pm$ 0.26 & 82.10 $\pm$ 0.64 & 92.77 $\pm$ 0.24 & \textbf{80.14 $\pm$ 1.07} & 87.76 \\
TIES-Merging & 0 & 94.27 $\pm$ 0.24 & 78.55 $\pm$ 0.92 & \textbf{94.08 $\pm$ 0.20} & 75.52 $\pm$ 1.25 & 85.61 \\
DARE & 0 & 96.03 $\pm$ 0.26 & 82.16 $\pm$ 0.68 & 92.81 $\pm$ 0.20 & \textbf{80.14 $\pm$ 0.87} & \textbf{87.78} \\
\textbf{NETA (Ours, $\alpha=1.0$)} & 0 & \textbf{96.29 $\pm$ 0.45} & \textbf{82.75 $\pm$ 0.80} & 92.61 $\pm$ 0.32 & 77.02 $\pm$ 0.56 & 87.17 \\
\midrule
\textit{Test-Time Optimization:} \\
Task-Wise AdaMerging & 4 & \textbf{98.49 $\pm$ 0.07} & 77.54 $\pm$ 0.00 & 89.70 $\pm$ 0.54 & 86.87 $\pm$ 1.32 & 88.15 \\
Layer-Wise AdaMerging & 52 & 98.44 $\pm$ 0.00 & \textbf{84.04 $\pm$ 1.02} & \textbf{92.92 $\pm$ 0.44} & \textbf{88.14 $\pm$ 1.31} & \textbf{90.89} \\
\bottomrule
\end{tabular}
\end{sc}
\end{small}
\end{center}
\vskip -0.1in
\end{table*}

\subsection{Quantitative Results and Discussion}
The primary experimental results are presented in Table~\ref{tab:results} and illustrated in Figure~\ref{fig:comparison}. We discuss our findings in detail below.

\begin{figure}[ht]
\vskip 0.2in
\begin{center}
\centerline{\includegraphics[width=\columnwidth]{comparison_plot.png}}
\caption{Visual comparison of multi-task model merging accuracy (\%) across MNIST, FashionMNIST, CIFAR-10, and SVHN. NETA provides robust, balanced performance without the catastrophic drop on FashionMNIST observed in Task-Wise AdaMerging.}
\label{fig:comparison}
\end{center}
\vskip -0.2in
\end{figure}

\subsubsection{Isotropic Magnitude Balancing (Ours vs. Task Arithmetic \& TIES-Merging)}
Standard Task Arithmetic directly merges task updates. However, during expert fine-tuning, tasks representing massive domain shifts—such as SVHN—undergo significantly larger parameter modifications compared to simpler datasets like MNIST or CIFAR-10. This leads to extremely high task vector norms for SVHN, which subsequently dominates the weight-space updates. Consequently, standard Task Arithmetic is heavily biased toward SVHN (80.14\% accuracy) while sacrificing performance on other tasks.

By equalizing the Frobenius norms of task vectors at every layer, \textbf{NETA} achieves isotropic representation strength across tasks. This closed-form balancing yields consistent improvements over standard Task Arithmetic and TIES-Merging on MNIST and FashionMNIST:
\begin{itemize}
    \item \textbf{MNIST}: NETA achieves \textbf{96.29\%} accuracy, outperforming Task Arithmetic (96.03\%, $+0.26\%$) and significantly outperforming TIES-Merging (94.27\%, $+2.02\%$).
    \item \textbf{FashionMNIST}: NETA achieves \textbf{82.75\%} accuracy, outperforming Task Arithmetic (82.10\%, $+0.65\%$) and outperforming TIES-Merging (78.55\%, $+4.20\%$).
\end{itemize}
These results demonstrate that enforcing a clean, physical constraint of magnitude equalization prevents high-norm updates from drowning out other tasks, establishing NETA as a highly competitive and often superior zero-shot alternative.

\textbf{Honest Trade-off and Average Performance.} On SVHN, NETA's accuracy is reduced (77.02\% vs. 80.14\% in Task Arithmetic), which causes NETA's average multi-task accuracy to be slightly lower than standard Task Arithmetic (87.17\% vs. 87.76\%). This behavior represents a classic \textit{peak performance vs. representation fairness trade-off}. In standard Task Arithmetic, the high-norm SVHN update hijacks the parameter space, achieving peak performance at the expense of others. NETA acts as an isotropic regularizer that prevents any single task from dominating. By equalizing the update norms, NETA deliberately curtails the dominance of SVHN to restore a balanced and fairer multi-task representation across the remaining three experts.

\subsubsection{The Overfitting-Optimizer Paradox (The AdaMerging Failure Mode)}
State-of-the-art Test-Time Adaptation (TTA) methods like AdaMerging minimize joint prediction entropy over calibration data to learn task and layer coefficients. Our results expose a severe vulnerability in this paradigm, which we designate as the \textbf{Overfitting-Optimizer Paradox}.

Under Task-Wise AdaMerging (optimizing 4 parameters), we observe a catastrophic performance collapse on \textbf{FashionMNIST}, with accuracy dropping from 82.10\% in Task Arithmetic to a dismal \textbf{77.54\%} (a drop of \textbf{-4.56\%}). A similar drop is observed on \textbf{CIFAR-10}, where accuracy falls by \textbf{-3.07\%} to 89.70\%.

This failure is directly driven by the joint entropy minimization objective. The entropy of predictions is highly dependent on task difficulty: simpler tasks like MNIST yield extremely sharp, low-entropy predictions quickly, while harder datasets like FashionMNIST and CIFAR-10 exhibit flatter, higher-entropy prediction distributions. To minimize the combined prediction entropy across tasks, the optimizer acts opportunistically. It maximizes the scaling parameters of simple, highly confident tasks (MNIST and SVHN) to near-maximum, while suppressing the scaling parameters of high-entropy, harder tasks (such as FashionMNIST and CIFAR-10) from the default $0.30$ down to approximately $0.23 - 0.24$. While these optimized scaling coefficients remain non-zero, this relative suppression is sufficient to degrade their weight representations and cause the observed performance collapse, resulting in a transductively overfit state where difficult tasks are actively suppressed and neglected.

Crucially, we clarify that the **Overfitting-Optimizer Paradox is a direct consequence of the unsupervised nature of the test-time adaptation objective** (i.e., joint prediction entropy minimization on unlabeled data). Because prediction entropy acts only as an unsupervised proxy for confidence rather than true accuracy, the optimizer is easily misled by task difficulty imbalances. Under a hypothetical *supervised* calibration objective (e.g., cross-entropy optimization on a small labeled validation set), this severe task suppression would likely not occur, as the objective function would explicitly penalize misclassifying the harder task. However, supervised adaptation is severely limited in practice, as it demands high-quality labeled downstream data, whereas NETA operates entirely zero-shot, data-free, and training-free.

\textbf{Statistical Notes and Boundary Convergence.} We also note a subtle empirical phenomenon in Table~\ref{tab:results}: Task-Wise AdaMerging on FashionMNIST and Layer-Wise AdaMerging on MNIST exhibit standard deviations of exactly \textbf{0.00\%} across the three independent random seeds. Through rigorous log verification, we confirm that this is not a checkpoint reuse artifact, but is instead driven by physical optimization boundaries. Specifically, under extreme task imbalances, the prediction entropy gradients consistently drive the corresponding scaling coefficients to their exact lower or upper parameter constraints (clamping boundaries) across all seeds. Furthermore, because evaluations are conducted on a discretized subset of 1024 test images, these stable clamped weights translate into identical classification counts, yielding a standard deviation of exactly $0.00\%$ while other unconstrained tasks retain standard non-zero variances. This high-dimensional clamping further exposes the lack of optimization stability in the test-time adaptation paradigm.

\textbf{Generality Boundaries of the Overfitting-Optimizer Paradox.} Why does the Overfitting-Optimizer Paradox catastrophically collapse Task-Wise AdaMerging but spare its Layer-Wise counterpart? The boundary conditions of this paradox lie in the \textit{spatial degrees of freedom} afforded to the optimization algorithm. 
In Task-Wise AdaMerging, the parameter space is severely constrained (only $K = 4$ scalars, one per task). Because the scaling factor is globally uniform across all layers, the optimizer is forced to make coarse, global adjustments: it must scale up a task's overall weight vector or globally suppress it. Under the joint prediction entropy minimization objective, simpler tasks with sharp, low-entropy predictions (like MNIST) yield massive gradients that dominate the optimization step. The optimizer opportunistically maximizes these simple tasks and, since it cannot adjust layers individually, it suppresses the overall coefficients of harder, flat-entropy tasks (like FashionMNIST and CIFAR-10) globally to reduce joint entropy. 

Conversely, Layer-Wise AdaMerging operates over $L \times K = 52$ continuous parameters, providing a substantial spatial buffer. With 13 parameter groups per task, the optimizer has the flexibility to selectively suppress and amplify task contributions at different depths of the network. For example, it can suppress a difficult task in early visual feature extraction layers (where representation sharing is high) while strongly amplifying it in deeper, task-specific layers, or vice-versa. This high-dimensional parameterization allows the optimizer to satisfy the joint entropy minimization objective without resorting to a global, severe suppression of any individual task, thereby recovering high test accuracy. 

However, we argue that this spatial flexibility is highly overparameterized, transductive, and prone to local calibration-set overfitting. Shuffling these delicate 52 coefficients post-optimization breaks this spatial coordination and collapses performance, highlighting that Layer-Wise AdaMerging achieves its peak performance (an average multi-task accuracy of \textbf{90.89\%}) at an unfavorable cost to engineering complexity and optimization stability.

NETA, by contrast, operates entirely zero-shot, data-free, and parameter-free. By mathematically enforcing isotropic magnitude balance, NETA avoids both the risk of catastrophic task suppression in low-dimensional spaces and the extreme complexity, data requirements, and transductive overfitting of high-dimensional test-time optimization loops. For deployments with strict compute limits, zero calibration data, or instant setup requirements, NETA offers an advantageous, parameter-free compromise that satisfies Occam's razor.

\subsection{Ablation Studies}
\label{sec:ablations}

To validate our architectural decisions and evaluate the practical formulations of NETA, we conduct a series of detailed ablation studies.

\begin{table}[ht]
\caption{Ablation study on CLIP ViT-B/32 merging (mean \% $\pm$ standard deviation across 3 seeds) for NETA configurations. We evaluate continuous $\alpha$-relaxation, the composite Group 0 grouping, the noise-damping stabilizer $\beta$, and the closed-form $\gamma^l$ update scale compensation.}
\label{tab:ablations}
\vskip 0.15in
\begin{center}
\begin{small}
\begin{sc}
\begin{tabular}{lccccr}
\toprule
Config & MNIST (\%) & FashionMNIST (\%) & CIFAR-10 (\%) & SVHN (\%) & Avg (\%) \\
\midrule
NETA ($\alpha=1.0$, default $\beta=10^{-6}$) & 96.29 $\pm$ 0.45 & 82.75 $\pm$ 0.80 & 92.61 $\pm$ 0.32 & 77.02 $\pm$ 0.56 & 87.17 \\
\textbf{NETA ($\alpha=0.5$)} & 96.16 $\pm$ 0.37 & 82.62 $\pm$ 0.60 & 92.71 $\pm$ 0.28 & \textbf{78.55 $\pm$ 0.65} & \textbf{87.51} \\
NETA (No Group 0) & 96.26 $\pm$ 0.41 & 82.71 $\pm$ 0.68 & \textbf{92.77 $\pm$ 0.24} & 76.99 $\pm$ 0.49 & 87.18 \\
\midrule
\textit{Noise-Damping Stabilizer $\beta$:} \\
NETA ($\beta = 10^{-3}$) & 96.29 $\pm$ 0.44 & 82.75 $\pm$ 0.80 & 92.61 $\pm$ 0.32 & 76.99 $\pm$ 0.56 & 87.16 \\
NETA ($\beta = 10^{-2}$) & 96.19 $\pm$ 0.44 & 82.68 $\pm$ 0.77 & 92.64 $\pm$ 0.36 & 76.76 $\pm$ 0.68 & 87.07 \\
\midrule
\textit{Scale Compensation $\gamma^l$:} \\
NETA + $\gamma^l$ Compensation & \textbf{96.32 $\pm$ 0.41} & \textbf{82.85 $\pm$ 0.77} & 92.61 $\pm$ 0.28 & 77.34 $\pm$ 0.57 & 87.28 \\
\bottomrule
\end{tabular}
\end{sc}
\end{small}
\end{center}
\vskip -0.1in
\end{table}

\textbf{Continuous $\alpha$-Relaxation.}
As presented in Equation~\ref{eq:alpha_relaxed} (Section~\ref{sec:method}), the continuous $\alpha$-relaxation parameter allows practitioners to smoothly interpolate between Task Arithmetic ($\alpha=0.0$) and full NETA ($\alpha=1.0$). 
In Table~\ref{tab:ablations}, we evaluate NETA under \textbf{$\alpha = 0.5$}. 
We find that $\alpha=0.5$ successfully resolves our peak performance trade-off:
\begin{itemize}
    \item On \textbf{SVHN}, NETA ($\alpha=0.5$) recovers performance to \textbf{78.55\%} (an improvement of \textbf{$+1.53\%$} compared to full NETA).
    \item Meanwhile, on \textbf{MNIST} (96.16\%) and \textbf{FashionMNIST} (82.62\%), it continues to outperform standard Task Arithmetic (96.03\% and 82.10\% respectively).
\end{itemize}
This continuous relaxation provides a highly flexible mechanism, letting practitioners optimize task-sharing dynamics based on domain-specific requirements entirely zero-shot.

\textbf{Composite Layer Grouping (Group 0).}
During initialization, standard NETA groups shallow embedding and non-transformer visual parameters into a composite Group 0 rather than normalizing them independently. 
In Table~\ref{tab:ablations}, we evaluate an ablation where these parameters are normalized independently (\textbf{No Group 0}). 
We find that omitting composite grouping degrades performance across MNIST (96.26\% vs 96.29\%) and FashionMNIST (82.71\% vs 82.75\%), while increasing the update scale on SVHN. This validates our physical heuristic: low-dimensional embedding parameters undergo minute fine-tuning shifts, and scaling them independently leads to representation noise and spatial distortion in early visual processing layers. Grouping them stabilizes the norm calculation and preserves downstream relative representation scale.

\textbf{Empirical Validation of the Noise-Damping Stabilizer ($\beta$).}
In Section~3.3, we theorized that setting the noise-damping stabilizer $\beta$ to a slightly larger, physically-grounded threshold (e.g., $10^{-3}$ or $10^{-2}$) acts as a soft-thresholding filter to prevent the unstable scaling of near-zero, noisy task vectors. In Table~\ref{tab:ablations}, we empirically evaluate this hypothesis across three independent seeds. 
We find that NETA is highly stable and robust to the choice of $\beta$. When increasing $\beta$ from the default $10^{-6}$ to $10^{-3}$, the classification accuracy remains virtually identical (achieving an average of $87.16\%$ compared to the default $87.17\%$, with negligible changes on MNIST and FashionMNIST). Even under a significantly larger stabilizer $\beta = 10^{-2}$, NETA maintains strong isotropic stabilization, achieving an average accuracy of $87.07\%$. This empirical validation confirms our theoretical claims that NETA's representation scaling is highly robust, and that practitioners can comfortably apply a noise-damping stabilizer to protect the stability of intermediate layers under negligible downstream updates.

\textbf{Empirical Evaluation of the Closed-Form Scale Compensation Factor ($\gamma^l$).}
To combat the directional norm contraction of the merged update vector analytically and training-free, we proposed a closed-form scale-compensation factor $\gamma^l$ in Equation~\ref{eq:compensation_factor} (Section~3.3). This factor rescales the merged NETA update at each layer to match the cumulative update norm of standard Task Arithmetic.
The empirical results of this closed-form extension are presented in Table~\ref{tab:ablations} (\textbf{NETA + $\gamma^l$ Compensation}). 
Applying $\gamma^l$ successfully yields a significant boost across all evaluation metrics, raising NETA's overall multi-task average accuracy to \textbf{87.28\%} ($+0.11\%$ improvement over standard NETA). Specifically, on MNIST and FashionMNIST, NETA + $\gamma^l$ achieves its highest accuracies of \textbf{96.32\%} and \textbf{82.85\%} respectively (improvements of $+0.29\%$ and $+0.75\%$ over standard Task Arithmetic). Crucially, on SVHN, applying $\gamma^l$ recovers performance from $77.02\%$ up to \textbf{77.34\%} ($+0.32\%$ absolute recovery), without requiring any manual grid search over $\lambda_0$. This empirical validation demonstrates that our proposed closed-form norm-restoration factor $\gamma^l$ represents an exceptionally powerful, training-free, and automated extension of NETA, successfully mitigating directional norm contraction entirely zero-shot.

\textbf{Representation Scale Analysis and $\lambda_0$ Grid Search.}
As theoretically qualified in Section~3.4, NETA's isotropic norm balancing contracts the overall merged update vector norm, which can introduce a slight performance deficit under a default scaling of $\lambda_0 = 0.30$. To analyze this geometric behavior, we conduct a systematic grid search over the global scaling coefficient $\lambda_0 \in [0.20, 0.40]$ on a representative 256-image subset. The quantitative results of this sweep are detailed in Table~\ref{tab:lambda_grid}.

The empirical findings presented in Table~\ref{tab:lambda_grid} demonstrate that standard Task Arithmetic is highly sensitive to the scale of $\lambda_0$, reaching its optimal average accuracy of $89.16\%$ at $\lambda_0 = 0.40$. When both methods are evaluated at the same default coefficient of $\lambda_0 = 0.30$, NETA's average accuracy undergoes a minor contraction (achieving $87.60\%$ for full NETA and $87.89\%$ for relaxed NETA, compared to $88.18\%$ for Task Arithmetic). If both methods are compared fairly across identical global coefficients, standard Task Arithmetic remains slightly superior in average accuracy across almost all columns of the grid. This confirms that NETA's isotropic normalization does not resolve update contraction to strictly outperform Task Arithmetic; rather, when both are tuned to their respective peak performance, standard Task Arithmetic maintains a marginal advantage ($89.16\%$ vs. $89.06\%$).

However, NETA's primary scientific value lies not in maximizing raw average accuracy via dominant updates, but in acting as an isotropic regularizer that enforces representation fairness. For example, at $\lambda_0 = 0.40$, relaxed NETA ($\alpha = 0.5$) improves FashionMNIST accuracy to $86.72\%$ compared to $85.94\%$ for standard Task Arithmetic, and full NETA ($\alpha = 1.0$) elevates it to $87.11\%$. This demonstrates that while standard Task Arithmetic achieves peak average performance by letting high-magnitude expert weights (such as SVHN) dominate the representation space, NETA deliberately curtails this dominance to distribute representation strength more equitably. This grid search proves that adjusting the global scaling coefficient $\lambda_0$ can effectively compensate for update contraction, bringing NETA's overall performance extremely close to optimal Task Arithmetic while successfully guaranteeing multi-task representation fairness.

\begin{table}[ht]
\caption{Global scaling coefficient $\lambda_0$ grid search results (Average Accuracy \% on 256-image subset). Across identical global coefficients, standard Task Arithmetic remains slightly superior in average accuracy, but NETA acts as an isotropic regularizer to guarantee multi-task representation fairness.}
\label{tab:lambda_grid}
\vskip 0.15in
\begin{center}
\begin{small}
\begin{sc}
\begin{tabular}{lcccccc}
\toprule
Method & $\lambda_0 = 0.20$ & $\lambda_0 = 0.25$ & $\lambda_0 = 0.30$ & $\lambda_0 = 0.33$ & $\lambda_0 = 0.35$ & $\lambda_0 = 0.40$ \\
\midrule
Task Arithmetic & 83.59 & 86.43 & 88.18 & 88.57 & 88.77 & 89.16 \\
NETA ($\alpha = 1.0$) & 83.11 & 85.94 & 87.60 & 88.28 & 88.48 & 88.48 \\
NETA ($\alpha = 0.5$) & 83.01 & 86.52 & 87.89 & 88.48 & 88.38 & 89.06 \\
\bottomrule
\end{tabular}
\end{sc}
\end{small}
\end{center}
\vskip -0.1in
\end{table}

\subsection{Complexity Analysis and Occam's Razor}
From a complexity and engineering perspective, NETA offers a highly favorable trade-off. It requires:
\begin{itemize}
    \item \textbf{Zero} optimization parameters (compared to 52 in Layer-Wise AdaMerging).
    \item \textbf{Zero} test-time calibration images (compared to 256 in AdaMerging).
    \item \textbf{Zero} training epochs, learning rates, or backpropagation passes.
\end{itemize}
Yet, NETA manages to achieve balanced, robust, and competitive multi-task accuracy entirely zero-shot, purely through a three-line analytical weight-space transform. This proves that a solid physical understanding of neural weight spaces can easily outperform parameter-heavy optimization pipelines, validating the power of simplicity in modern deep learning.
